% mixfrag.tex — The TKF92 Model with Fragment Mixtures (MixFrag)
\section{The TKF92 Model with Fragment Mixtures (MixFrag)}
\label{sec:mixfrag}

MixFrag generalises TKF92 (Section~\ref{sec:tkf92}) by promoting the single
fragment-extension parameter $\ext$ to a categorical latent variable.
When a mortal link's fragment is created, it is independently assigned a
\emph{fragtype} $\frag \in \{1,\ldots,\nfrag\}$ with
$\frag \sim \catdist(\fragdist_1,\ldots,\fragdist_\nfrag)$
(weights $\fragdist_\frag \geq 0$, $\sum_\frag \fragdist_\frag = 1$),
and the fragment then consists of $\fraglen \sim \geomdist(\ext_\frag)$
characters from $\alphabet$, all governed by the same substitution process
$\subproc(\exch,\eqm)$.
In the notation of Section~\ref{sec:tkf92},
\begin{multline*}
\seq'(\evoltime) \sim
\mixfrag\!\left(\subproc(\exch,\eqm);\insrate,\delrate,
\{\ext_\frag\},\{\fragdist_\frag\}\right) \\
\equiv
\tkflinks\!\left(\fragproc^\nfrag\!\left(\subproc(\exch,\eqm);
\{\ext_\frag\},\{\fragdist_\frag\}\right);\insrate,\delrate\right)
\end{multline*}
where the fragment process first draws a fragtype and then a fragtype-specific
geometric run,
\begin{multline*}
\state \sim \fragproc^\nfrag(\model;\{\ext_\frag\},\{\fragdist_\frag\})
\ \Leftrightarrow\
\frag \sim \catdist(\fragdist_1,\ldots,\fragdist_\nfrag), \\
\fraglen \sim \geomdist(\ext_\frag),\
\state = (x_1 \ldots x_\fraglen),\ x_\sumidx \sim \model.
\end{multline*}
The only parameters beyond TKF91 are the $\nfrag$ fragtype-specific extension
probabilities $\ext_\frag$ and the $\nfrag$ fragtype weights $\fragdist_\frag$;
the indel (BDI) process and the substitution process are both shared across
fragtypes.
With $\nfrag = 1$ (so $\fragdist_1 = 1$) MixFrag reduces to TKF92.
It uses a single substitution model and a single, non-nested indel process, and
so introduces no null states.


\subsection{Singlet HMM and Pair HMM Representations}
\label{sec:mixfrag-machines}

We do not give a WFST for MixFrag.
The fragtype is latent information that is not visible in the sequence, so a
transducer that respects fragment indivisibility over evolutionary time must
expose the fragtype in its interface alphabet---in order to keep
branch-to-branch composition consistent---making the construction even more
elaborate than the TKF92 WFST of Appendix~\ref{sec:tkf92-wfst}.
The Singlet and Pair HMMs, by contrast, follow systematically from TKF91,
exactly as for TKF92.

The same construction yields both machines from their TKF91 counterparts: a
transition \emph{out of} an emitting state of fragtype $\srcfrag$ acquires a
$(1-\ext_\srcfrag)$ multiplier (the current fragment terminates); a transition
\emph{into} an emitting state of fragtype $\destfrag$ acquires a
$\fragdist_\destfrag$ multiplier (the new fragment is assigned fragtype
$\destfrag$); and each emitting state acquires an added self-loop term
$\ext_\srcfrag\,\delta_{\srcfrag\destfrag}$ that extends the current fragment
within its own fragtype.

\paragraph{Singlet HMM (stationary distribution).}
The Singlet HMM has $2+\nfrag$ states $\{\sta,\ins_\frag,\fin\}$.
Applying the construction to the TKF91 stationary HMM (which emits
$n \sim \geomdist(\kappa)$ characters, $\kappa = \insrate/\delrate$) gives the
transition matrix
\[
\left(
\begin{array}{r|ccc}
& \sta & \ins_\destfrag & \fin \\
\hline
\sta          & 0 & \kappa\,\fragdist_\destfrag & 1-\kappa \\
\ins_\srcfrag & 0 & \ext_\srcfrag\,\delta_{\srcfrag\destfrag}
                    + (1-\ext_\srcfrag)\,\kappa\,\fragdist_\destfrag
                  & (1-\ext_\srcfrag)(1-\kappa) \\
\fin          & 0 & 0 & 0
\end{array}
\right)
\]
where $\ins_\srcfrag$ emits a character from $\eqm$.
Each $\ins_\srcfrag$ row sums to
$\ext_\srcfrag + (1-\ext_\srcfrag)\big(\kappa\sum_\destfrag\fragdist_\destfrag
+ (1-\kappa)\big) = \ext_\srcfrag + (1-\ext_\srcfrag) = 1$,
using $\sum_\destfrag\fragdist_\destfrag = 1$.

\paragraph{Pair HMM (finite-time joint distribution).}
The Pair HMM has $2+3\nfrag$ states
$\{\sta,\mat_\frag,\ins_\frag,\del_\frag,\fin\}$;
for convenience we write $\sta_\frag \equiv \sta$ and $\fin_\frag \equiv \fin$,
so that any state may be written $\xstate_\frag$ with
$\xstate \in \{\sta,\mat,\ins,\del,\fin\}$.
Writing $\tkftrans_{\xstate\ystate} \equiv
\tkftrans_{\xstate\ystate}(\insrate,\delrate,\evoltime)$ for the TKF91 Pair HMM
transition probabilities (Section~\ref{sec:tkf91-machines}), the joint Pair HMM
transition matrix is
\[
\resizebox{\textwidth}{!}{$\displaystyle
\tkftrans^\nfrag(\insrate,\delrate,\evoltime,\{\ext_\frag\},\{\fragdist_\frag\}) =
\left(
\begin{array}{r|ccccc}
& \sta & \mat_\destfrag & \ins_\destfrag & \del_\destfrag & \fin \\
\hline
\sta & 0
  & \fragdist_\destfrag\tkftrans_{\sta\mat}
  & \fragdist_\destfrag\tkftrans_{\sta\ins}
  & \fragdist_\destfrag\tkftrans_{\sta\del}
  & \tkftrans_{\sta\fin} \\
\mat_\srcfrag & 0
  & \ext_\srcfrag\delta_{\srcfrag\destfrag} + (1-\ext_\srcfrag)\fragdist_\destfrag\tkftrans_{\mat\mat}
  & (1-\ext_\srcfrag)\fragdist_\destfrag\tkftrans_{\mat\ins}
  & (1-\ext_\srcfrag)\fragdist_\destfrag\tkftrans_{\mat\del}
  & (1-\ext_\srcfrag)\tkftrans_{\mat\fin} \\
\ins_\srcfrag & 0
  & (1-\ext_\srcfrag)\fragdist_\destfrag\tkftrans_{\ins\mat}
  & \ext_\srcfrag\delta_{\srcfrag\destfrag} + (1-\ext_\srcfrag)\fragdist_\destfrag\tkftrans_{\ins\ins}
  & (1-\ext_\srcfrag)\fragdist_\destfrag\tkftrans_{\ins\del}
  & (1-\ext_\srcfrag)\tkftrans_{\ins\fin} \\
\del_\srcfrag & 0
  & (1-\ext_\srcfrag)\fragdist_\destfrag\tkftrans_{\del\mat}
  & (1-\ext_\srcfrag)\fragdist_\destfrag\tkftrans_{\del\ins}
  & \ext_\srcfrag\delta_{\srcfrag\destfrag} + (1-\ext_\srcfrag)\fragdist_\destfrag\tkftrans_{\del\del}
  & (1-\ext_\srcfrag)\tkftrans_{\del\fin} \\
\fin & 0 & 0 & 0 & 0 & 0
\end{array}
\right)
$}
\]
Emissions are inherited unchanged from TKF92 and are independent of fragtype:
$\mat_\frag$ aligns an ancestral--descendant residue pair under
$\exp(\revsub\evoltime)$, while $\ins_\frag$ and $\del_\frag$ emit a single
inserted or deleted residue from $\eqm$.
Each emitting row sums to 1: for $\xstate \in \{\mat,\ins,\del\}$, since
$\sum_\destfrag\fragdist_\destfrag = 1$ the fragtype weights collapse and the row
sum reduces to
$\ext_\srcfrag + (1-\ext_\srcfrag)\sum_\ystate \tkftrans_{\xstate\ystate}
= \ext_\srcfrag + (1-\ext_\srcfrag) = 1$,
exactly as for TKF92.
The Pair HMM has no null states, so no null-elimination (collapse) step is
required.


\subsection{Baum-Welch Algorithm for MixFrag}
\label{sec:bw-mixfrag}

The Baum-Welch algorithm follows that for TKF92 (Section~\ref{sec:bw-tkf92}), but
the fragment-extension and fragment-end statistics are now accumulated per
fragtype.

\paragraph{E-step.}
Run Forward-Backward on the MixFrag Pair HMM, yielding expected transition counts
$\hat{n}'_{a_\srcfrag b_\destfrag}$ over the $2+3\nfrag$ states.
For an emitting state $a_\frag$ ($a \in \{\mat,\ins,\del\}$), the self-loop count
$\hat{n}'_{a_\frag a_\frag}$ combines fragment extension within fragtype $\frag$
(probability $\ext_\frag$) with a new same-state link that is again assigned
fragtype $\frag$ (probability $(1-\ext_\frag)\tkftrans_{aa}\fragdist_\frag$).
We resolve this self-loop and recover the TKF91 link-level counts, accumulating
the per-fragtype fragment-extension count $F_\frag$ and fragment-end count
$E_\frag$ (the latter is also the expected number of type-$\frag$ fragments,
since each fragment terminates exactly once):
\begin{eqnarray*}
F_{a\frag} & = & \frac{\hat{n}'_{a_\frag a_\frag}\,\ext_\frag}
                      {\ext_\frag + (1-\ext_\frag)\tkftrans_{aa}\fragdist_\frag},
                 \quad a \in \{\mat,\ins,\del\} \\
F_\frag & = & \sum_{a \in \{\mat,\ins,\del\}} F_{a\frag} \\
E_\frag & = & \sum_{a \in \{\mat,\ins,\del\}} \sum_{b_\destfrag}
              \hat{n}'_{a_\frag\, b_\destfrag} \ -\ F_\frag \\
\hat{n}_{a_\srcfrag b_\destfrag} & = & \hat{n}'_{a_\srcfrag b_\destfrag}
              - \delta_{ab}\,\delta_{\srcfrag\destfrag}\,F_{a\srcfrag}
\end{eqnarray*}
The denominator of $F_{a\frag}$ is the full weight of the $a_\frag \to a_\frag$
self-loop: its fragment-extension branch ($\ext_\frag$) and its
new-link-same-state branch (which must re-draw the same fragtype, hence the
$\fragdist_\frag$ factor).
Summing the resolved counts over fragtypes,
$\hat{n}_{ab} = \sum_{\srcfrag,\destfrag} \hat{n}_{a_\srcfrag b_\destfrag}$,
recovers the TKF91 link-level transition counts, from which $S,B,D,W,U,L,M,V$ are
accumulated exactly as for TKF91 (the indel and substitution processes being
shared across fragtypes); in particular
$L = \sum_{i}(\hat{n}_{i\mat} + \hat{n}_{i\del})$ counts links (fragments of any
type) and the character-level statistics $W,U,V$ are accumulated per character as
in TKF91.
Because the fragtype is latent, the Forward-Backward E-step marginalises it
together with the alignment, so the per-fragtype counts $F_\frag, E_\frag$ are
posterior expectations rather than counts read directly off a given alignment.

\paragraph{M-step.}
Identical to TKF91, except that the fragtype-specific parameters are set from the
per-fragtype fragment statistics,
\[
\ext_\frag \leftarrow \frac{F_\frag}{F_\frag + E_\frag},
\qquad
\fragdist_\frag \leftarrow \frac{E_\frag}{\sum_{\frag'} E_{\frag'}}.
\]
Here $E_\frag$ is the expected number of fragments of fragtype $\frag$ (each
fragment terminates exactly once), so $\fragdist_\frag$ is the categorical MLE
over fragtypes, while $\ext_\frag$ is the geometric MLE, giving mean fragment
length $1/(1-\ext_\frag) = (F_\frag + E_\frag)/E_\frag$ for fragtype $\frag$.
